\section{CARGO-AnoT}

\subsection{Evidence interface over a normal rule graph}

For a candidate fact $x$, AnoT produces a base score $a(x)$ and a set of projected and reachable rules. \method converts these rule traversals into a named evidence vector $\phi(x)$. The generic portion contains rule coverage, support, missing-precursor, temporal-gap, successor-conflict, relation-pair conflict, and subject-state conflict features. Each feature is computed by deterministic graph traversal in \texttt{rule\_graph\_features.py}. For example, the support feature is the fraction of expected predecessor rules that are observed; the missing-precursor feature is its complement. A temporal-gap feature compares the observed gap with historical gap statistics, and successor conflict measures disagreement between projected rules and their reachable successors.

For the explicit-fusion paths, the final score combines the AnoT score and active evidence features:
\begin{equation}
S_k(x)=\alpha_k a(x)+\sum_{j\in\mathcal{F}_k}w_{k,j}\,\psi_{k,j}(x),
\label{eq:fusion}
\end{equation}
where $k$ denotes the error type, $\mathcal{F}_k$ is its feature set, and $\psi_{k,j}$ applies the type-specific sign convention. In particular, support is negative evidence for ordinary anomaly scoring but positive evidence for a missing fact. The evaluator selects weights and a decision threshold on validation and freezes them before test evaluation. For the conceptual \texttt{fair\_select} and optional missing calibration paths, the same named features are instead passed to a validation-selected linear, logistic, or shallow gradient-boosted calibrator. Thus, Equation~\ref{eq:fusion} describes the explicit path rather than every implementation branch.

Table~\ref{tab:evidence-interface} makes the interface explicit. The base graph and generic features are shared infrastructure, not separate detectors. The type-specific features are attached only to the matching task. This restriction prevents a feature with the wrong direction from being imported merely because it improved a different anomaly category on validation.

\begin{table}[t]
\centering
\caption{Evidence interface used by the full configuration. ``Positive direction'' means that a larger feature value increases the score for the stated task after applying $\psi_{k,j}$.}
\label{tab:evidence-interface}
\small
\begin{tabular}{@{}p{0.23\linewidth}p{0.39\linewidth}p{0.22\linewidth}@{}}
\toprule
Task & Active evidence & Positive direction \\
\midrule
Conceptual & Pair novelty; subject--relation and object--relation novelty; entity--relation novelty; role support; neighborhood overlap & Role incompatibility / novelty \\
Time & Base temporal features; validated contrastive atomic-gap rules & Conflict / displaced-gap match \\
Missing & Context coverage; target support; pair and entity--relation history; temporal support & Support for the held-out target \\
\bottomrule
\end{tabular}
\end{table}

The explicit fusion implementation applies the sign before weighting. For instance, conceptual relation-role support receives a negative sign because repeated role evidence makes a conceptual corruption less likely. In contrast, missing support receives a positive sign because a well-supported but absent target is the positive class in that task. The evaluator serializes every component value and the final score for each example; calibration metadata records the selected model when a learned path is used.

\subsection{Conceptual evidence from relational roles}

Entity-pair novelty alone is unreliable in event data because valid new pairs appear frequently. The conceptual module therefore builds train-graph profiles for subject--relation and object--relation counts, relation-specific pair history, and entity-neighborhood overlap. Given $(s,r,o,t)$, it derives features such as whether $s$ has previously occurred as the subject of $r$, whether $o$ has occurred as the object of $r$, and the normalized support for these roles. The implementation also includes low-neighbor-overlap and pair-relation novelty features.

The module can mine type-relation incompatibility rules from training perturbations, but the reported conceptual configuration uses a validation-selected \texttt{novelty\_role} feature set. The \texttt{fair\_select} implementation compares linear, logistic, and shallow gradient-boosted candidates on a temporal holdout within validation; its gradient-boosted candidates use $80$ estimators, depth $1$--$3$, learning rates $0.02$, $0.05$, or $0.1$, and random state $42$. No test labels take part in feature selection or calibration.

The role profile is intentionally finer than pair novelty. A previously unseen pair $(s,o)$ is common in event streams and is not, by itself, a reliable error signal. The useful question is whether $s$ has occurred in the subject role of $r$, whether $o$ has occurred in its object role, and whether their local relation neighborhoods are compatible with the observed role. The selected \texttt{novelty\_role} set contains pair novelty, pair--relation novelty, generated-sample novelty, subject--relation novelty, object--relation novelty, entity--relation novelty, relation-role support, and low-neighbor-overlap features. The reported time configuration uses explicit score fusion; missing support can use either its explicit path or its separately validated optional calibration path.

\subsection{Contrastive temporal-gap rules}

Time errors require a different signal. We construct training references that keep a fact's subject, relation, and object but place it at a later target time. The generator takes facts from the earliest ten training timestamps, requires that the fact's latest training occurrence precede the target by at least $200$ time units, and excludes facts known in validation or test. At each target time, it samples up to $30\%$ as many displaced references as normal facts. With ICEWS14s, this yields $25{,}794$ normal training samples and $7{,}654$ displaced references.

For each sample, the miner extracts a pattern
\begin{equation}
q=(r_{\mathrm{pre}},r_{\mathrm{tar}},[\ell,u)),
\end{equation}
where $r_{\mathrm{pre}}$ and $r_{\mathrm{tar}}$ are oriented atomic rules on the same entity pair and $[\ell,u)$ is a gap bucket. The default buckets are $[0,30)$, $[30,90)$, $[90,180)$, $[180,365)$, and $[365,\infty)$. Repeated-rule patterns are allowed: an old occurrence of the same atomic rule can be informative when it reappears at an implausible time.

For each pattern, we count unique normal and displaced training samples, calculate a label-partition MDL gain, and require a positive anomaly direction. Candidates must satisfy at least $10$ anomaly hits, $20$ total hits, lift at least $1.25$, AUC at least $0.505$, posterior direction probability at least $0.90$, a positive $95\%$ credible lower bound on the rate difference, and positive MDL gain. Validation recomputes the statistics on normal facts and AnoT time errors; it does not create new patterns. The selected-rule weight combines normalized training MDL gain, validation enrichment, and posterior direction probability.

More precisely, for a candidate $q$, let $h^+_q,n^+_q$ be its anomaly hits and total anomaly samples, and let $h^-_q,n^-_q$ be the corresponding normal quantities. The validation selector uses
\begin{equation}
\operatorname{AUC}_q=\frac{1}{2}\left(\frac{h^+_q}{n^+_q}+1-\frac{h^-_q}{n^-_q}\right),
\quad
\operatorname{lift}_q=
\frac{h^+_q/(h^+_q+h^-_q)}{n^+_q/(n^+_q+n^-_q)}.
\label{eq:rule-statistics}
\end{equation}
It additionally draws from independent Beta--Binomial posteriors for the two hit rates and requires $\Pr(p_q^+>p_q^-)\geq0.90$ together with a positive lower endpoint of the $95\%$ credible interval for $p_q^+-p_q^-$. The final bounded weight is proportional to
\begin{equation}
\widetilde{\operatorname{MDL}}_q\,\max\{0,\log_2(\operatorname{lift}_q)\}\,
\Pr(p_q^+>p_q^-),
\label{eq:rule-weight}
\end{equation}
where $\widetilde{\operatorname{MDL}}_q$ is the training MDL gain normalized among eligible rules. Equations~\ref{eq:rule-statistics}--\ref{eq:rule-weight} explain why a frequent but nondiscriminative temporal pattern is discarded rather than becoming a conflict rule.

\subsection{Discovery, validation, and freezing}

The temporal rule file is constructed in three stages. In discovery, the miner creates normal/displaced training batches and extracts all observed atomic-gap patterns. It first removes patterns that fail minimum hit counts, reverse the anomaly direction, have weak lift or AUC, or do not reduce the label-partition description length after paying the rule-encoding cost. In validation, it replays the \emph{fixed} candidate list on validation normal and time-error samples. This second stage computes the posterior tests in Equation~\ref{eq:rule-statistics}, ranks surviving rules by posterior direction, credible rate difference, AUC, and anomaly hits, and retains at most $500$ rules. Finally, it writes a serialized rule file. Test evaluation loads only this file; it never enumerates a new temporal pattern.

The distinction between discovery and validation is more than an implementation detail. A time-displaced reference is a controlled contrast source, not a claim that every displaced event is a real anomaly. The train stage uses it to propose hypotheses about typed contexts. The validation stage asks whether the same hypotheses separate the benchmark's independently generated time errors from normal facts. A candidate can be statistically prominent in the train reference data and still be rejected on validation. Conversely, a rule's eventual test contribution is traceable to a bounded candidate set selected before test scoring.

This procedure selected $500$ temporal-gap rules from $1{,}255$ candidates. Their union matched $65.4\%$ of validation time errors and $3.2\%$ of validation normal samples, with union AUC $0.811$. One representative pattern links an out-oriented \texttt{Make\_statement} atomic rule to the same rule with a $[180,365)$ gap. It does not assert that repeated statements are anomalous. It records that this typed, oriented, long-gap recurrence is much more frequent in the displaced reference context than in normal data.

\subsection{Train-only support for missing facts}

Missing-error evidence has the opposite polarity. The missing-support miner uses only training facts and represents a rule as $(r_{\mathrm{ctx}},r_{\mathrm{tar}},\rho,[\ell,u))$, where $r_{\mathrm{ctx}}$ is a contextual atomic rule, $r_{\mathrm{tar}}$ is a target rule, $\rho$ specifies a shared subject, object, or pair role, and $[\ell,u)$ is a gap bucket. It retains patterns with sufficient context and target support, then exposes contextual coverage, temporal support, pair history, entity-relation support, and broad support features. The primary missing configuration uses this train-only support graph; the optional contrastive missing-support variant is disabled because its validation behavior was not stable.

\subsection{Auditable candidate workflow and explanations}

The evidence workflow wraps the detector without replacing its decision rule. Candidate artifacts include source, scope, polarity, binding, gap interval, dataset hash, and provenance. CandidateAgent can package symbolic, contrastive, conceptual, missing-support, or optional LLM-proposed candidates. The reported time-rule audit uses deterministic candidate sources. EvidenceAgent retrieves matched historical facts and rule evidence. ValidationAgent accepts a candidate only when its validation evidence passes the configured enrichment, AUC, Beta--Binomial direction, credible-difference, and MDL checks. ScoringAgent converts accepted candidates into the rule JSON consumed by the CARGO-AnoT evaluation path.

Every agent inherits a common trace writer that records input and output hashes, tool calls, validation status, runtime, and failure reason. ExplanationAgent produces a natural-language explanation only from accepted evidence artifacts. ExplanationVerifierAgent then checks that the target fact, historical fact, rule identifier, and time gap in the explanation are all present in the evidence record. The workflow therefore treats an explanation as a checked view of a score trace, not as a separate generator of reasons.

The workflow has eight fixed responsibilities. CandidateAgent emits a versioned candidate specification; EvidenceAgent retrieves historical matches; ValidationAgent applies split-specific tests; ScoringAgent writes the frozen rule artifact; ExplanationAgent renders an accepted artifact for a user; ExplanationVerifierAgent checks every referenced field; EvaluationAgent records scores and metrics; and ReflectionAgent aggregates failure reasons for subsequent development. Only the first four stages can change the rule file, and only before test evaluation. The later stages may inspect artifacts but cannot add a rule or update the graph. This separation is important in an online setting, where an explanation or a post-hoc diagnosis must not create an unlogged feedback path into the detector.

The workflow is deliberately typed rather than prompt-driven. A candidate specifies its source, error polarity, rule binding, temporal interval, dataset scope, and provenance. An evidence record lists the candidate identifier, target fact, matched historical fact, observed interval, and label/split context. A validation record holds the computed statistics and explicit rejection reasons. The scoring artifact contains only accepted rules and their bounded weights. This contract lets the code reject malformed or split-inconsistent artifacts before they influence a score. It also leaves room for a future LLM candidate proposer without making an LLM part of the present detection claim: any such proposal would still need the same evidence retrieval and deterministic validation.

